% !Mode:: "TeX:UTF-8"
% 正文三章，照范例的段落序列；章号 1 / 4 / 6、节号跳着来也照它

\chapter{绪论}[Introduction]

\section{课题背景及研究的目的和意义}[Background, objective and significance of the subject]

发展国防工业、微电子工业等尖端技术需要精密和超精密的仪器设备，精密仪器设备要求高速、……

……

\section{气体润滑轴承及其相关理论的发展概况}[Developmental of gas-lubricated bearing and correlated theories]

气体轴承是利用气膜支撑负荷或减少摩擦的机械构件。……

……

\subsection{气体润滑轴承的发展}[Developmental of gas-lubricated bearing]

1828年，R.R.Willis\cite{zhao1998}发表了一篇关于小孔节流平板中压力分布的文章，这是有记载的研究气体润滑的最早文献。……

根据间隙内气膜压力的产生原理，气体轴承可以分为四种基本形式：

（1）气体静压轴承\quad 加压气体经过节流器进入间隙，在间隙内产生压力气膜使物体浮起的气体轴承，……

\subsection{气体润滑轴承的分类}[Classification of gas-lubricated bearing]

根据间隙内气膜压力的产生原理，气体轴承可以分为四种基本形式，其结构如图\ref{fig:1-1}所示。

（1）气体静压轴承\quad 加压气体经过节流器进入间隙，在间隙内产生压力气膜使物体浮起的气体轴承，结构如图\ref{fig:1-1} (a)所示。……

……

\begin{figure}[!ht]
  \centering
  \begin{minipage}{\textwidth}
    \centering
    \subfigure[]{\includegraphics[width=6.92cm]{aerostatic}}
    \hspace{1cm}
    \subfigure[]{\includegraphics[width=6.65cm]{aerodynamic}}
  \end{minipage}
  \begin{minipage}{\textwidth}
    \centering
    \subfigure[]{\includegraphics[width=6.72cm]{hybrid}}%
    \hspace{0.5cm}%
    \subfigure[]{\includegraphics[width=7.71cm]{squeeze-film}}
  \end{minipage}
  \caption{气体润滑轴承的分类}[Classification of the gas-lubricated bearing]
  \label{fig:1-1}

  {\wuhao (a) 气体静压轴承; (b) 气体动压轴承; (c) 气体动静压轴承; (d) 气体压膜轴承\par}
\end{figure}

（也可以按照下图范例书写）

% 这一张是图 1-1 换个写法的示范，范例里仍编 1-1；博士那份带题注，本科那份没有
% （\sampleAltFigureCaption 各是各的）。续图接着编 (c)(d)，靠 \ContinuedFloat。
% 子图之间的行尾要用 % 吃掉空格：6.70 + 1 + 7.17 = 14.87 cm，再加两个词间空格就
% 超过 15 cm 的版心，两张子图折成上下两行。上面那张第二排 6.72 + 7.71 已经 14.43，
% 只留得下半厘米的空。
\begin{figure}[!ht]
  \centering
  \subfigure[气体静压轴承]{\includegraphics[width=6.87cm]{aerostatic}}%
  \hspace{1cm}%
  \subfigure[气体动压轴承]{\includegraphics[width=6.03cm]{aerodynamic}}
  \sampleAltFigureCaption
\end{figure}

\clearpage

\begin{figure}[!ht]
  \ContinuedFloat
  \centering
  \subfigure[气体动静压轴承]{\includegraphics[width=6.70cm]{hybrid}}%
  \hspace{1cm}%
  \subfigure[气体压膜轴承]{\includegraphics[width=7.17cm]{squeeze-film}}
  \caption*{图\thefigure（续图）}
\end{figure}

\setcounter{subsection}{4}
\subsection{多孔质气体静压轴承的研究}[Research on porous externally pressurized gas bearing]

由于气体的压力低和可压缩性，……。

\subsubsection{多孔质静压轴承的分类}[Classification of porous thrust bearing]

轴承工作面的整体或……。

\subsubsection{多孔质材料特性的研究}[Research on characteristics of porous materials]

材料的主要特点是具有一定的……。

（1）孔隙特性\quad 多孔质材料是由……。

……

\setcounter{section}{3}
\section{本文的主要研究内容}[Main research content of this subject]

本课题的研究内容主要是针对局部多孔质止推轴承的多孔质材料的渗透率、静压轴承的静态特性、稳定性及其影响因素进行展开，……。

\setcounter{chapter}{3}
\chapter{基于FLUENT软件的轴承静态特性研究}[Research on static characteristics of bearing based on FLUENT software]

\section{引言}[Introduction]

利用现成的商用软件来研究流场，可以免去对N-S方程求解程序的……

\setcounter{section}{2}
\setcounter{subsection}{2}
\subsection{边界条件的设定}[Initialization of boundary conditions]

本文采用……，则每一个方向上的……由公式\eqref{eq:4-1}、\eqref{eq:4-2}求得：
\begin{equation}\label{eq:4-1}
  \alpha = \frac{D_{\mathrm{p}}^2}{150}\frac{\psi^3}{(1-\psi)^2}
\end{equation}
\begin{equation}\label{eq:4-2}
  C_2 = \frac{3.5}{D_{\mathrm{p}}}\frac{(1-\psi)}{\psi^3}
\end{equation}
\begin{eqdenote}
  \item $D_{\mathrm{p}}$ & 多孔质材料的平均粒子直径（m）；
  \item $\psi$ & 孔隙度（孔隙体积占总体积的百分比）；
  \item $C_2$ & 特征渗透性或固有渗透性（m$^2$）。
\end{eqdenote}

……

\setcounter{section}{3}
\section{本章小结}[Brief summary]

……

\setcounter{chapter}{5}
\chapter{局部多孔质静压轴承的试验研究}[Experiment on partial porous thrust bearing]

\section{引言}[Introduction]

在前面几章中，分别对局部多孔质材料的渗透率……

\section{多孔质石墨渗透率测试试验}[Experiments on permeability of porous graphite]

……

1号试样的试验数据见表\ref{tab:6-1}。

\begin{table}[!ht]
  \centering
  % 范例里英文题注的温度那一截是另起一行的（原件是两段），照它断行
  \caption{1号试样渗透率测试数据(温度：$T$=16 ℃\quad 高度：$H$=5.31 mm)}[Data of measured permeability of sample 1\\ (Temperature: $T$=16 ℃\quad Height: $H$=5.31 mm)]
  \label{tab:6-1}
  % 范例的格子是 vAlign=center，表头那几格行数不一样，居中才与原件同位
  \begin{tblr}{colspec={X[c]X[c]X[c]X[c]X[c]X[c]}, width=\textwidth, rows={m}}
    \toprule
    {供气压力\\ $P_{\mathrm{s}}$ (MPa)} & {流量测量\\ $M'$ (m$^3$/h)} & {流量修正值 $M$ (m$^3$/s)\\ $\times$10$^{-4}$} & {压力差\\ $\Delta P$ (Pa)} & lg$\Delta P$ & lg $M$ \\
    \midrule
    0.15 & 0.009 & 0.023 12 & 46 900 & 4.671 17 & -5.636 01 \\
    0.2 & 0.021 & 0.045 84 & 96 900 & 4.986 32 & -5.338 76 \\
    0.25 & 0.039 & 0.074 13 & 146 900 & 5.167 02 & -5.130 01 \\
    0.3 & 0.097 & 0.167 47 & 196 900 & 5.294 24 & -4.776 06 \\
    0.35 & 0.136 & 0.217 53 & 246 900 & 5.392 52 & -4.662 48 \\
    0.4 & 0.171 & 0.254 85 & 296 900 & 5.472 61 & -4.593 72 \\
    0.45 & 0.202 & 0.284 67 & 346 900 & 5.540 20 & — \\
    \bottomrule
  \end{tblr}
\end{table}

% 范例在这行前留了一串空段（博士那份十三个、本科那份十二个，数在主文件里；修正版
% 按正文格式，1.25 倍不贴格），换页落在哪由空段自己排：博士那份 Word 排到下一页
% 顶上还剩三个，本科那份全在这一页装下
\hitenter*[\sampleEntersAfterTable]

……如果表格超过一页时，可采用续表的形式移入下一页：

\addtocounter{table}{-1}
% 修正版范例的续表只有中文表题，没有英文那一行
\begin{longtblr}[
  caption={试样渗透率测试数据},
  label={tab:6-1b}]{colspec={X[c]X[c]X[c]X[c]X[c]X[c]}, width=\textwidth, rows={m}}
  \toprule
  {供气压力\\ $P_{\mathrm{s}}$ (MPa)} & {流量测量\\ $M'$ (m$^3$/h)} & {流量修正值 $M$ (m$^3$/s)\\ $\times$10$^{-4}$} & {压力差\\ $\Delta P$ (Pa)} & lg$\Delta P$ & lg $M$ \\
  \midrule
  0.15 & 0.009 & 0.023 12 & 46 900 & 4.671 17 & -5.636 01 \\
  0.2 & 0.021 & 0.045 84 & 96 900 & 4.986 32 & -5.338 76 \\
  0.25 & 0.039 & 0.074 13 & 146 900 & 5.167 02 & -5.130 01 \\
  0.15 & 0.009 & 0.023 12 & 46 900 & 4.671 17 & -5.636 01 \\
  0.2 & 0.021 & 0.045 84 & 96 900 & 4.986 32 & -5.338 76 \\
  0.25 & 0.039 & 0.074 13 & 146 900 & 5.167 02 & -5.130 01 \\
  0.15 & 0.009 & 0.023 12 & 46 900 & 4.671 17 & -5.636 01 \\
  0.2 & 0.021 & 0.045 84 & 96 900 & 4.986 32 & -5.338 76 \\
  0.25 & 0.039 & 0.074 13 & 146 900 & 5.167 02 & -5.130 01 \\
  0.15 & 0.009 & 0.023 12 & 46 900 & 4.671 17 & -5.636 01 \\
  0.2 & 0.021 & 0.045 84 & 96 900 & 4.986 32 & -5.338 76 \\
  0.25 & 0.039 & 0.074 13 & 146 900 & 5.167 02 & -5.130 01 \\
  0.15 & 0.009 & 0.023 12 & 46 900 & 4.671 17 & -5.636 01 \\
  0.2 & 0.021 & 0.045 84 & 96 900 & 4.986 32 & -5.338 76 \\
  0.25 & 0.039 & 0.074 13 & 146 900 & 5.167 02 & -5.130 01 \\
  0.15 & 0.009 & 0.023 12 & 46 900 & 4.671 17 & -5.636 01 \\
  0.2 & 0.021 & 0.045 84 & 96 900 & 4.986 32 & -5.338 76 \\
  0.25 & 0.039 & 0.074 13 & 146 900 & 5.167 02 & -5.130 01 \\
  $\vdots$ & & & & & \\
  0.3 & 0.097 & 0.167 47 & 196 900 & 5.294 24 & -4.776 06 \\
  0.35 & 0.136 & 0.217 53 & 246 900 & 5.392 52 & -4.662 48 \\
  0.4 & 0.171 & 0.254 85 & 296 900 & 5.472 61 & -4.593 72 \\
  0.45 & 0.202 & 0.284 67 & 346 900 & 5.540 20 & — \\
  \bottomrule
\end{longtblr}

\setcounter{section}{4}
\section{本章小结}[Brief summary]

……
